\documentclass[10pt,a4paper]{amsart}
\usepackage[margin=19mm]{geometry}
\usepackage{amsmath,amssymb,mathtools}
\usepackage[scr=rsfs,cal=euler]{mathalfa}
\usepackage{xfrac}
\usepackage[unicode,hidelinks,bookmarks=false]{hyperref}
\newcommand{\ie}{i.e.\ }
\newcommand{\cf}{cf.\ }
\newcommand{\resp}{respectively\ }
\newcommand{\Subsection}{\subsection}
\providecommand{\nobreakdash}{\nobreak\hbox{-}}
\setlength{\emergencystretch}{2em}
\multlinegap0pt
\usepackage{amssymb}
\usepackage{mathtools}
\usepackage{xcolor}
\usepackage[shortlabels]{enumitem}
\newtheorem{proposition}{Proposition}
\renewcommand{\Re}{\operatorname{Re}}
\newcommand{\eps}{\varepsilon}
\title{\Huge \textbf{Regular Arithmetic Functions}\\ \vspace{1.4em} \huge \textsc{Volume I}\\ \vspace{0.7em} \Large Theory, Applications, Examples}
\author{Benoit Cloitre}
\date{Source: arXiv:2609.09366v1 (CC BY 4.0)}
\begin{document}
\maketitle

\noindent\emph{Source and licence.} \Huge \textbf{Regular Arithmetic Functions}\\ \vspace{1.4em} \huge \textsc{Volume I}\\ \vspace{0.7em} \Large Theory, Applications, Examples. Version arXiv:2609.09366v1 (CC BY 4.0), distributed by arXiv under the Creative Commons Attribution 4.0 International licence, http://creativecommons.org/licenses/by/4.0/, as recorded in the arXiv licence metadata for this exact version and shown on the abstract page https://arxiv.org/abs/2609.09366v1. Full licence notice: \texttt{licenses/arxiv-2609.09366-CC-BY-4.0.txt}.\par\smallskip

\noindent\textit{Editorial scope.} Counted unit: the article's proposition thm:affine\_quadratic\_gauge (source0.tex lines 10342--10442). The article's complete original proof of it is delivered with it (source0.tex lines 10444--10637, one \texttt{proof} environment of 194 lines). No mathematics is written, repaired or re-derived: the statement and the proof are the article's own lines, in the article's own notation, and no retained line is substituted or re-flowed.  No further source-local context is needed: every cross-reference of every retained line resolves inside the two delivered spans.  Nothing was omitted from any retained span, so the package carries no omission record.  No retained line contains a citation, so the wrapper carries no bibliography and no bibliography entry.  No retained line contains a figure environment, an image-inclusion command or drawing code, so no image is delivered and the package declares no asset dependency.  Editorial formatting only, in addition: an AMS \texttt{amsart} preamble that restates the article's own declarations and macros used by the retained lines, namely the environments \texttt{proposition} on separate counters, together with the standard packages those lines need.  The retained environments print in this document as Proposition 1 in order of appearance, whereas the article prints the same items with the numbering of its own full document.  The complete native source of the article as distributed in the arXiv e-print for this exact version is bundled unchanged as \texttt{sources/209809/source0.tex}.
\begin{proposition}\label{thm:affine_quadratic_gauge}
For each real \(\beta\), the gauged equation has a unique solution. Put
\[
A(n)=\sum_{k\leq n}a_k,
\qquad
B(n)=\sum_{k\leq n}k^2a_k.
\]
Then \(A(1)=a_1=1\), and for every \(n\geq2\),
\[
A(n)=c_nA(n-1)+d_n,
\]
where
\[
c_n=(1-\lambda)+\lambda\left(1-\frac1n\right)^2
=1-\frac{2\lambda}{n}+\frac{\lambda}{n^2}
\]
and
\[
d_n=
\frac{n^{2-2\beta}-(n-1)^{2-2\beta}}{n^2}.
\]
With \(P_1=1\) and
\[
P_n=\prod_{j=2}^{n}c_j,
\]
there is a constant \(C_\lambda>0\) such that
\[
P_n=C_\lambda n^{-2\lambda}
\left(1+\mathcal O(n^{-1})\right).
\]
The exact variation formula is
\[
A(n)=P_n\left(
1+\sum_{m=2}^{n}\frac{d_m}{P_m}
\right).
\]

The integral defining the transform of the composed kernel converges
initially for \(\Re s<0\) and gives
\[
h^*(s)
=-s\int_0^1h(t)t^{-s-1}\,dt
=\frac{2\lambda-s}{2-s}.
\]
Its meromorphic continuation has a simple zero at \(s=2\lambda\) and a
simple pole at \(s=2\). Define the transform in the quadratic gauge from
this kernel-side integral by
\[
G^*_f(z):=h^*(2z),
\qquad \Re z<0.
\]
It has the meromorphic continuation
\[
G^*_f(z)
=\frac{\lambda-z}{1-z}
=g^*(z).
\]
Its simple zero at \(z=\lambda\) does not cancel with its simple pole at
\(z=1\).

The partial sums have the following three behaviors. If \(\beta<\lambda\),
then
\[
A(n)\sim
\frac{1-\beta}{\lambda-\beta}n^{-2\beta}
=\frac{f(n)^{-\beta}}{G^*_f(\beta)}.
\]
If \(\beta=\lambda\), then
\[
A(n)\sim
2(1-\lambda)n^{-2\lambda}\log n.
\]
If \(\beta>\lambda\), then
\[
A(n)=\mathcal O(n^{-2\lambda})
=\mathcal O\!\left(f(n)^{-\lambda}\right).
\]
For \(\beta=1\), more precisely,
\[
d_n=0
\quad(n\geq2),
\qquad
A(n)=P_n\sim C_\lambda n^{-2\lambda}.
\]

The sharp ordinary transition point of the composed kernel and the sharp
gauged transition point of the original kernel are
\[
\alpha(h)=2\lambda,
\qquad
\alpha_f(G)=\frac{\alpha(h)}2=\lambda.
\]
Moreover
\[
\eta(h)=2\lambda,
\qquad
\eta_f(G)=\lambda.
\]
Thus the affine kernel is stable and in equilibrium under the quadratic
gauge.
\end{proposition}

\begin{proof}
Since \(g(1)=1\), the diagonal coefficient of the triangular equation is
one. The solution is unique, and the equation at \(n=1\) gives
\[
a_1=A(1)
=\frac{f(1)^{-\beta}}{G(f(1),f(1))}
=1.
\]
Multiplication of the gauged equation by \(n^2\) gives the exact identity
\[
\lambda n^2A(n)+(1-\lambda)B(n)=n^{2-2\beta}.
\]
The increments of the two sums satisfy
\[
B(n)-B(n-1)
=n^2a_n
=n^2\bigl(A(n)-A(n-1)\bigr).
\]
Insert this identity into the equation at \(n\), and use the equation at
\(n-1\) to eliminate \(B(n-1)\). The result is
\[
n^2A(n)
=\left((1-\lambda)n^2+\lambda(n-1)^2\right)A(n-1)
+n^{2-2\beta}-(n-1)^{2-2\beta}.
\]
Division by \(n^2\) proves the stated recurrence without any asymptotic
input.

The transform is also constructed before any asymptotic conclusion. The
constant term of \(h\) shows that its integral converges exactly in the
initial half plane \(\Re s<0\). There
\[
\begin{aligned}
h^*(s)
&=-s\left(
(1-\lambda)\int_0^1t^{1-s}\,dt
+\lambda\int_0^1t^{-s-1}\,dt
\right)\\
&=-s\left(
\frac{1-\lambda}{2-s}-\frac{\lambda}{s}
\right)\\
&=\frac{2\lambda-s}{2-s}.
\end{aligned}
\]
This rational expression supplies the meromorphic continuation. For
\(\Re z<0\), the substitution \(x=t^2\) gives
\[
\begin{aligned}
G^*_f(z)
&=-2z\int_0^1g(t^2)t^{-2z-1}\,dt\\
&=-z\int_0^1g(x)x^{-z-1}\,dx\\
&=g^*(z)
=\frac{\lambda-z}{1-z}.
\end{aligned}
\]
The inequalities \(0<\lambda<1\) separate the zero \(z=\lambda\) from the
pole \(z=1\). At this point the transform identifies only the analytic
candidate \(\eta_f(G)=\lambda\). It does not identify the arithmetic index.

It remains to obtain that index from the recurrence. For \(j\geq2\), set
\[
u_j=c_j-1=-\frac{2\lambda}{j}+\frac{\lambda}{j^2}.
\]
The numbers \(c_j\) lie in \((0,1)\), and Taylor expansion\index[terms]{Taylor expansion}\index[names]{Taylor, B.} of the logarithm
with a uniform remainder gives
\[
\log c_j=-\frac{2\lambda}{j}+\mathcal O(j^{-2}).
\]
Consequently the series
\[
\sum_{j=2}^{\infty}
\left(\log c_j+\frac{2\lambda}{j}\right)
\]
converges. The harmonic-sum expansion and the tail estimate for this series
show that a finite real number \(L_\lambda\) satisfies
\[
\log P_n=-2\lambda\log n+L_\lambda+\mathcal O(n^{-1}).
\]
Exponentiation yields
\[
P_n=C_\lambda n^{-2\lambda}
\left(1+\mathcal O(n^{-1})\right),
\qquad
C_\lambda=e^{L_\lambda}>0.
\]

Since \(P_n=c_nP_{n-1}\), division of the recurrence by \(P_n\) gives
\[
\frac{A(n)}{P_n}
=\frac{A(n-1)}{P_{n-1}}+\frac{d_n}{P_n}.
\]
Summation from \(2\) to \(n\), together with \(A(1)=P_1=1\), proves the
exact variation formula.

Suppose first that \(\beta\ne1\). The elementary power-difference estimate
gives
\[
d_m
=2(1-\beta)m^{-2\beta-1}
\left(1+\mathcal O(m^{-1})\right).
\]
It follows that
\[
\frac{d_m}{P_m}
=\frac{2(1-\beta)}{C_\lambda}
m^{2(\lambda-\beta)-1}
\left(1+\mathcal O(m^{-1})\right).
\]

If \(\beta<\lambda\), then \(\beta<1\), and the power sum diverges with
\[
\sum_{m=2}^{n}\frac{d_m}{P_m}
\sim
\frac{1-\beta}{C_\lambda(\lambda-\beta)}
n^{2(\lambda-\beta)}.
\]
Multiplication by \(P_n\) gives
\[
A(n)\sim
\frac{1-\beta}{\lambda-\beta}n^{-2\beta}.
\]
Since
\[
G^*_f(\beta)=\frac{\lambda-\beta}{1-\beta},
\]
this is exactly the stated transparency formula.

If \(\beta=\lambda\), then
\[
\frac{d_m}{P_m}
\sim
\frac{2(1-\lambda)}{C_\lambda}\frac1m.
\]
The harmonic sum gives
\[
A(n)\sim2(1-\lambda)n^{-2\lambda}\log n.
\]

If \(\beta>\lambda\) and \(\beta\ne1\), the series
\[
\sum_{m=2}^{\infty}\left|\frac{d_m}{P_m}\right|
\]
converges because its summand is
\(
\mathcal O\bigl(m^{-1-2(\beta-\lambda)}\bigr)
\).
The variation formula then gives
\[
A(n)=\mathcal O(P_n)=\mathcal O(n^{-2\lambda}).
\]
In fact, if
\[
L_{\lambda,\beta}
=1+\sum_{m=2}^{\infty}\frac{d_m}{P_m},
\]
then
\[
n^{2\lambda}A(n)\longrightarrow
C_\lambda L_{\lambda,\beta}.
\]
When \(\beta=1\), the numerator defining \(d_n\) is identically zero for
every \(n\geq2\). Hence \(A(n)=P_n\), which proves both the required bound
and the stated nonzero equivalent.

The same calculation can now be read in the ordinary exponent
\(\gamma=2\beta\). It gives transparency for every \(\gamma<2\lambda\),
logarithmic resonance at \(\gamma=2\lambda\), and homogeneous decay of order
\(n^{-2\lambda}\) above that value. Thus the sharp ordinary transition point
is
\[
\alpha(h)=2\lambda.
\]
The passage to the gauged index is quantitative. For \(\gamma=2\beta\),
\[
\frac{n^{-\gamma}}{h^*(\gamma)}
=\frac{f(n)^{-\beta}}{G^*_f(\beta)}.
\]
For the absorption estimate, choosing \(\delta=2\eps\) gives
\[
\mathcal O(n^{-2\lambda+\delta})
=\mathcal O\!\left(f(n)^{-\lambda+\eps}\right).
\]
The transparency formula holds up to \(\lambda\) from below and fails at
\(\lambda\) because of the logarithmic factor. The case \(\beta=1\) gives
the nonzero homogeneous equivalent
\(A(n)\sim C_\lambda f(n)^{-\lambda}\), so the absorption exponent cannot
be improved. Therefore \(\lambda\) is the sharp gauged transition point and
\[
\alpha_f(G)=\frac{\alpha(h)}2=\lambda.
\]
The transform calculation independently gives \(\eta_f(G)=\lambda\). Since
the undeformed theorem gives \(\alpha(G)=\lambda\), stability and equilibrium
follow.
\end{proof}

\end{document}
